% ============================================================
\chapter{Couverture Delta et Gestion des Risques}
\label{ch:couverture}
% ============================================================

Le mod\`ele de Black-Scholes ne se contente pas de donner un prix~: il dit comment \emph{se couvrir}. L'id\'ee centrale est que le risque d'une option peut \^etre \'elimin\'e en d\'etenant en permanence la bonne quantit\'e du sous-jacent --- le \emph{delta}. En ajustant cette position \`a chaque instant, le portefeuille devient insensible aux fluctuations du march\'e. Mais la couverture parfaite est un id\'eal~: en pratique, on ne peut r\'eajuster qu'\`a intervalles discrets, les co\^uts de transaction s'accumulent, et la volatilit\'e n'est jamais constante. Comprendre ces limites est aussi important que comprendre la th\'eorie.

\section{Introduction}

La couverture (\textit{hedging}) est l'activit\'e consistant \`a r\'eduire
ou \'eliminer le risque li\'e \`a une position sur produits d\'eriv\'es.
Ce chapitre d\'eveloppe la th\'eorie et la pratique du delta-hedging,
du gamma-hedging, et analyse les co\^uts et limites de la couverture
dynamique.

\section{Delta-hedging}

\subsection{Principe}

\begin{definition}[Delta-hedging]
Le \textbf{delta-hedging} consiste \`a maintenir un portefeuille
$\Pi_t = V_t - \Delta_t S_t$
localement insensible aux variations de $S$, o\`u
$\Delta_t = \frac{\partial V}{\partial S}(S_t, t)$.
\end{definition}

\begin{theoreme}[Portefeuille auto-finan\c{c}ant delta-couvert]
Dans le cadre de Black--Scholes avec trading continu, le portefeuille
delta-couvert est sans risque et r\'eplique parfaitement le d\'eriv\'e.
Sa valeur satisfait~:
\[
    d\Pi_t = r\Pi_t\,dt.
\]
\end{theoreme}

\begin{intuition}[title=\textbf{Couverture en temps discret}]
En pratique, le rebalancement est discret (toutes les $\Delta t$ unit\'es
de temps). L'erreur de couverture provient du fait que $\Delta$ change
entre deux rebalancements. Plus la fr\'equence est \'elev\'ee, meilleure
est la couverture, mais plus les co\^uts de transaction augmentent.
\end{intuition}

\subsection{Erreur de couverture en temps discret}

\begin{proposition}[Erreur de couverture]
Pour un call europ\'een couvert en delta avec rebalancement aux dates
$t_0, t_1, \ldots, t_N$, l'erreur de couverture est~:
\[
    \epsilon = \sum_{i=0}^{N-1} \frac{1}{2}\Gamma(S_{t_i}, t_i)
    S_{t_i}^2\bigl[(\Delta S_{t_i})^2 - \sigma^2 S_{t_i}^2 \Delta t\bigr]
    \cdot e^{r(T-t_{i+1})},
\]
o\`u $\Delta S_{t_i} = S_{t_{i+1}} - S_{t_i}$.
\end{proposition}

L'erreur de couverture a une esp\'erance nulle sous $\mathbb{Q}$ mais
une variance non nulle, proportionnelle \`a $\Delta t$.

\begin{codeblock}[title=\textbf{Simulation du delta-hedging discret}]
\begin{minted}{python}
import numpy as np
from scipy.stats import norm

def bs_delta(S, K, T, r, sigma):
    """Delta d'un call europeen."""
    if T < 1e-10:
        return 1.0 if S > K else 0.0
    d1 = (np.log(S/K) + (r + 0.5*sigma**2)*T) / (sigma*np.sqrt(T))
    return norm.cdf(d1)

def bs_price(S, K, T, r, sigma):
    """Prix Black-Scholes d'un call."""
    if T < 1e-10:
        return max(S - K, 0)
    d1 = (np.log(S/K) + (r + 0.5*sigma**2)*T) / (sigma*np.sqrt(T))
    d2 = d1 - sigma*np.sqrt(T)
    return S*norm.cdf(d1) - K*np.exp(-r*T)*norm.cdf(d2)

def simulate_delta_hedge(S0, K, T, r, sigma, N_rebal, M=10000):
    """Simule le delta-hedging et calcule l'erreur de couverture."""
    dt = T / N_rebal
    errors = np.zeros(M)

    for m in range(M):
        S = S0
        # Portefeuille initial: vente du call, achat de delta*S
        V0 = bs_price(S, K, T, r, sigma)
        delta = bs_delta(S, K, T, r, sigma)
        cash = V0 - delta * S  # position cash

        for i in range(N_rebal):
            tau = T - (i+1)*dt
            Z = np.random.standard_normal()
            S_new = S * np.exp((r - 0.5*sigma**2)*dt
                               + sigma*np.sqrt(dt)*Z)
            # Cash croît au taux sans risque
            cash *= np.exp(r * dt)
            # Nouveau delta
            if tau > 1e-10:
                delta_new = bs_delta(S_new, K, tau, r, sigma)
            else:
                delta_new = 1.0 if S_new > K else 0.0
            # Rebalancement
            cash -= (delta_new - delta) * S_new
            delta = delta_new
            S = S_new

        # Valeur finale du portefeuille
        portfolio = cash + delta * S
        payoff = max(S - K, 0)
        errors[m] = portfolio - payoff

    return errors

np.random.seed(42)
for N in [10, 50, 252, 1000]:
    errs = simulate_delta_hedge(100, 100, 1.0, 0.05, 0.20, N)
    print(f"N={N:4d}: mean={np.mean(errs):.4f}, "
          f"std={np.std(errs):.4f}")
\end{minted}
\end{codeblock}

\section{Gamma-hedging}

\begin{definition}[Gamma-hedging]
Le \textbf{gamma-hedging} consiste \`a rendre le portefeuille insensible
non seulement aux petites variations de $S$ (delta-neutre) mais aussi
aux variations plus importantes (gamma-neutre). On utilise un deuxi\`eme
instrument (typiquement une autre option) pour annuler le gamma.
\end{definition}

\begin{proposition}[Construction d'un portefeuille gamma-neutre]
Soit $V$ le d\'eriv\'e \`a couvrir et $O$ un instrument de couverture
avec gamma $\Gamma_O$. La position dans $O$ n\'ecessaire pour annuler
le gamma est~:
\[
    w_O = -\frac{\Gamma_V}{\Gamma_O}.
\]
Apr\`es cette couverture, on ajuste la position dans le sous-jacent pour
restaurer la delta-neutralit\'e~:
\[
    \Delta_{\text{total}} = \Delta_V + w_O \Delta_O.
\]
\end{proposition}

\begin{exemple}
Un trader est short un call ($\Delta = 0.55$, $\Gamma = 0.03$). Il dispose
d'un call de maturit\'e diff\'erente ($\Delta_O = 0.40$, $\Gamma_O = 0.05$).

Position gamma-neutre~: $w_O = -(-0.03)/0.05 = 0.60$.

Nouveau delta~: $-0.55 + 0.60 \times 0.40 = -0.31$.

Il faut acheter $0.31$ unit\'es du sous-jacent pour \^etre delta-neutre.
\end{exemple}

\section{Vega-hedging}

\begin{definition}[Vega-hedging]
Le \textbf{vega-hedging} vise \`a neutraliser la sensibilit\'e du
portefeuille \`a la volatilit\'e. Comme le vega du sous-jacent est nul,
il faut utiliser des options.
\end{definition}

\begin{remarque}
En pratique, le vega-hedging est plus important que le gamma-hedging
car les mouvements de volatilit\'e implicite peuvent engendrer des
profits et pertes (P\&L) consid\'erables.
\end{remarque}

\section{D\'ecomposition du P\&L}

\begin{theoreme}[D\'ecomposition du P\&L d'un portefeuille delta-couvert]
Le P\&L instantan\'e d'un portefeuille delta-couvert s'\'ecrit~:
\[
    d\text{P\&L} = \frac{1}{2}\Gamma S^2\bigl[(dS/S)^2 - \sigma^2\,dt\bigr]
                 + \mathcal{V}\,d\sigma_{\text{impl}} + \cdots
\]
Le premier terme est le \textbf{P\&L gamma} (profit si la volatilit\'e
r\'ealis\'ee d\'epasse la volatilit\'e implicite), le second le
\textbf{P\&L vega}.
\end{theoreme}

\begin{formules}[title=\textbf{Approximation Taylor du P\&L}]
Pour un portefeuille g\'en\'eral, l'approximation Taylor donne~:
\[
    \Delta\text{P\&L} \approx \Delta\cdot\delta S
    + \frac{1}{2}\Gamma(\delta S)^2 + \Theta\,\delta t
    + \mathcal{V}\,\delta\sigma + \rho\,\delta r.
\]
\end{formules}

\section{Co\^uts de transaction et couverture optimale}

\begin{theoreme}[Fr\'equence de rebalancement optimale (Leland)]
En pr\'esence de co\^uts de transaction proportionnels $k\abs{\delta S}$,
Leland (1985) montre que la couverture optimale utilise une volatilit\'e
modifi\'ee~:
\[
    \hat{\sigma}^2 = \sigma^2\left(1 + \sqrt{\frac{2}{\pi}}
    \frac{k}{\sigma\sqrt{\Delta t}}\right).
\]
\end{theoreme}

\begin{proposition}[Mod\`ele de Whalley--Wilmott]
Une approche alternative consiste \`a ne rebalancer que lorsque le delta
s'\'ecarte suffisamment de sa cible. La bande de non-rebalancement
optimale est~:
\[
    \Delta_{\text{cible}} \pm \left(\frac{3k\,e^{-r\tau}\Gamma^2 S^2}
    {2\lambda}\right)^{1/3},
\]
o\`u $\lambda$ repr\'esente l'aversion au risque du trader.
\end{proposition}

\begin{codeblock}[title=\textbf{Couverture avec co\^uts de transaction}]
\begin{minted}{python}
def hedge_with_costs(S0, K, T, r, sigma, N_rebal, k=0.001, M=5000):
    """Delta-hedging avec couts de transaction proportionnels."""
    dt = T / N_rebal
    total_costs = np.zeros(M)
    hedge_errors = np.zeros(M)

    for m in range(M):
        S = S0
        V0 = bs_price(S, K, T, r, sigma)
        delta = bs_delta(S, K, T, r, sigma)
        cash = V0 - delta * S
        cost = k * abs(delta) * S  # cout initial

        for i in range(N_rebal):
            tau = T - (i+1)*dt
            Z = np.random.standard_normal()
            S_new = S * np.exp((r-0.5*sigma**2)*dt
                               + sigma*np.sqrt(dt)*Z)
            cash *= np.exp(r*dt)

            delta_new = bs_delta(S_new, K, max(tau, 1e-10), r, sigma)
            trade = delta_new - delta
            cost += k * abs(trade) * S_new
            cash -= trade * S_new
            delta = delta_new
            S = S_new

        portfolio = cash + delta * S
        payoff = max(S - K, 0)
        hedge_errors[m] = portfolio - payoff
        total_costs[m] = cost

    print(f"Cout moyen: {np.mean(total_costs):.4f}")
    print(f"Erreur std: {np.std(hedge_errors):.4f}")
    return hedge_errors, total_costs

np.random.seed(42)
hedge_with_costs(100, 100, 1.0, 0.05, 0.20, 252, k=0.001)
\end{minted}
\end{codeblock}

\section{Couverture en march\'e incomplet}

\begin{remarque}
Lorsque le march\'e est incomplet (volatilit\'e stochastique, sauts),
la couverture parfaite est impossible. On recourt alors \`a~:
\begin{itemize}
    \item La \textbf{couverture en variance minimale}~: minimiser
          $\Var(\Pi_T - H)$ sur les strat\'egies admissibles.
    \item La \textbf{couverture par super-r\'eplication}~: trouver le
          co\^ut minimal d'un portefeuille dominant le payoff.
    \item La \textbf{couverture par quantile}~: contr\^oler la probabilit\'e
          de perte.
\end{itemize}
\end{remarque}

\section{Exercices}

\begin{exercice}[Simulation du delta-hedging]
\begin{enumerate}
    \item Simuler la couverture delta d'un call europ\'een pour
          $N = 10, 50, 252, 1000$ rebalancements.
    \item Tracer l'histogramme de l'erreur de couverture pour chaque $N$.
    \item V\'erifier que la variance de l'erreur d\'ecro\^it en $O(1/N)$.
\end{enumerate}
\end{exercice}

\begin{exercice}[Gamma-hedging]
\begin{enumerate}
    \item Construire un portefeuille delta-gamma neutre \`a partir d'un
          call ATM court, d'un call OTM long et du sous-jacent.
    \item Simuler le P\&L sur une semaine et comparer avec le portefeuille
          uniquement delta-couvert.
\end{enumerate}
\end{exercice}

\begin{exercice}[Analyse du P\&L]
\begin{enumerate}
    \item D\'ecomposer le P\&L journalier d'un portefeuille delta-couvert
          en composantes delta, gamma, theta et vega.
    \item Montrer que pour un portefeuille delta-neutre,
          $\text{P\&L} \approx \frac{1}{2}\Gamma S^2(\sigma_r^2 - \sigma_i^2)\,dt$
          o\`u $\sigma_r$ est la volatilit\'e r\'ealis\'ee et $\sigma_i$
          la volatilit\'e implicite.
\end{enumerate}
\end{exercice}

\begin{exercice}[Co\^uts de transaction]
\begin{enumerate}
    \item Impl\'ementer la couverture de Leland avec volatilit\'e modifi\'ee.
    \item Comparer le co\^ut total et l'erreur de couverture avec la
          m\'ethode standard pour diff\'erentes valeurs de $k$.
    \item Impl\'ementer la bande de Whalley--Wilmott et analyser le
          compromis co\^ut/risque.
\end{enumerate}
\end{exercice}
